\headline{{\textbf{\Large\sffamily Seasonal Care Tips:}}\\
          {\textbf{\large\sffamily Mid\--February to Mid\--March}}}
\label{2025-02-care}

As the days grow longer and the chill of winter begins to loosen its grip, mid\--February to mid\--March is a pivotal period in bonsai care.
While spring isn't in full swing just yet, these weeks offer the perfect opportunity to prepare your bonsai for the upcoming growth season.
Whether you're a seasoned enthusiast or new to the art, this guide will help you navigate the late winter to early spring transition with confidence.

\medskip
\topic{A Weather\--Watchful Eye}
\noindent
In London's temperate climate, this time of year can be rather unpredictable.
Frosty mornings may still surprise us, even as we start to notice milder afternoons.
Your bonsai are also feeling these changes, and their care needs to adapt accordingly.

Ensure your trees remain protected from any lingering frosts.
If you've been overwintering them in a cold frame, greenhouse, or sheltered spot, keep an eye on overnight temperatures.
On milder days, it's a good idea to give them some fresh air and sunlight.
For those keeping indoor bonsai, make sure they're not in direct line of chilly draughts or too close to heat sources, as this can cause stress.

\medskip
\topic{Watering Wisely}
\noindent
One of the most common pitfalls during this period is overwatering.
While the soil may not dry out as quickly as in warmer months, it's crucial to monitor moisture levels.
The root systems are just beginning to awaken, and excess water can lead to root rot.

Check the soil with your finger.
If it feels dry a centimetre below the surface, it's time to water.
Use tepid water to avoid shocking the roots.
Remember, balance is key: neither a parched bonsai nor a soggy one will thrive.

\medskip
\topic{Pruning and Cleaning}
\noindent
Mid\--February to mid\--March is an excellent time for structural pruning on deciduous species, as they are still mostly dormant and bare of leaves.
This allows you to see the branch structure clearly and make thoughtful cuts to shape your tree.

Before making any cuts, clean and sharpen your tools.
Dull blades can tear the wood, leaving your bonsai vulnerable to disease.
Remove any deadwood, crossing branches, or overly vigorous shoots that disrupt the tree's balance.
For evergreens like junipers or pines, avoid heavy pruning now --- their growth cycle will soon pick up, and over\--pruning can weaken them.

Don't forget to tidy up the pot!
Remove moss, algae, or weeds from the soil surface.
These can compete with your bonsai for nutrients and water, so keeping the soil clean sets the stage for healthy growth.

\medskip
\topic{Repotting: Is It Time?}
\noindent
This is the prime time to repot many bonsai species, particularly deciduous trees and some evergreens.
As the roots begin to stir from dormancy, they are better equipped to recover from the stress of repotting.
However, timing is crucial --- wait until the buds are just starting to swell but haven't fully opened.

When repotting, inspect the roots carefully.
Trim back any circling or dead roots and ensure the remaining roots are neatly spread out.
Use a fresh, well\--draining bonsai soil mix suitable for your tree's species.
Repotting not only gives the roots room to grow but also improves aeration and nutrient uptake.

If your tree doesn't require repotting this year, consider refreshing the top layer of soil to boost its nutrient content.
Use a gentle rake or chopstick to remove the old soil and add a fresh layer of bonsai substrate.

\medskip
\topic{Feeding: Hold Your Horses}
\noindent
It's tempting to dive into feeding your bonsai as soon as you see the first signs of growth, but patience is key.
Overfeeding too early can cause a flush of growth that the tree may struggle to support.
Wait until the buds have opened and growth is underway before introducing fertiliser.

When it's time, start with a mild, balanced fertiliser or one with a slightly higher nitrogen content to encourage leaf and shoot growth.
For now, focus on healthy soil conditions and water management.

\medskip
\topic{Pest Patrol}
\noindent
As temperatures rise, so too does pest activity.
Late winter and early spring are the perfect time to inspect your bonsai for any overwintering pests or fungal issues.
Aphids, spider mites, and scale insects can lurk in the nooks and crannies of bark or branches.

Use a soft brush or cloth to clean the bark and inspect thoroughly.
If you spot any pests, address them promptly with a gentle insecticidal soap or neem oil.
Catching problems early prevents infestations from taking hold as the growing season begins.

\medskip
\topic{Planning Ahead}
\noindent
This period is not only about tending to the immediate needs of your bonsai but also about planning ahead.
Think about the goals you have for each tree this year.
Are you focusing on refining the shape, encouraging specific growth, or preparing for a show?
Adjust your care routine accordingly, and keep a bonsai journal to track progress.

\medskip
\topic{A Word on Patience}
\noindent
As bonsai enthusiasts, we know that patience is one of the most valuable virtues.
Mid\--February to mid\--March is a time of quiet preparation, setting the stage for the vibrant growth that lies ahead.
Each small step you take now contributes to the health and beauty of your bonsai in the months to come.

So, enjoy these transitional weeks!
Whether it's a cup of tea in hand as you check on your trees or the quiet satisfaction of a freshly repotted bonsai, this is a wonderful time to reconnect with the art of bonsai.

Happy bonsai tending!
\closearticle
\clearpage
